\documentclass[10pt,a4paper]{article}

\input{config}

% Paquetes adicionales usados en este main
\usepackage{longtable}
\usepackage{array}
\usepackage{booktabs}

% Marca visible para información pendiente
\newcommand{\falta}[1]{\textbf{\textcolor{red}{FALTA: #1}}}

\begin{document}

\input{portada}

%==============================================
% Índice
%==============================================
\pagenumbering{gobble}
\tableofcontents
\newpage
\pagenumbering{arabic}

%==============================================
% Desarrollo
%==============================================


\section{PYME y problema}
%==============================================

\subsection{Identificación de la PYME}

El proyecto se desarrollará para un minimarket de tamaño pequeño ubicado en Rengo, Región de O'Higgins. El negocio funciona actualmente con una gestión principalmente manual y visual del inventario y con un seguimiento financiero interno limitado.

\begin{itemize}[leftmargin=*]
    \item \textbf{Nombre de la PYME:} Alma Neira Muñoz
    \item \textbf{Rubro:} minimarket / comercio minorista de productos de consumo habitual.
    \item \textbf{Tamaño aproximado:} {Micro empresa (clasificación oficial SII), categoría tributaria Primera, afecta a IVA}
    \item \textbf{Ubicación:}Pje Carlos Isamitt 785 Rengo, Región de O'Higgins.
\end{itemize}

\subsection{Evidencia de existencia}

{La existencia de la PYME se verifica mediante Declaración Jurada de Inicio de Actividades emitida por el Servicio de Impuestos Internos (Folio N° 21007629601, trámite efectuado el 30-09-2021), que certifica el giro comercial (almacén, comestibles, bebidas, lácteos, frutas y verduras) y clasifica al negocio como micro empresa. El documento completo se adjunta en el Anexo A.}

\subsection{Problema de gestión}

El principal problema identificado es la ausencia de un sistema formal de gestión de inventario. Actualmente no existe un registro permanente que permita conocer con precisión cuánto stock hay de cada producto. La reposición se realiza principalmente de forma visual: cuando el encargado observa que queda poco de un producto, decide solicitar nuevas unidades al proveedor.

Esta forma de operación dificulta conocer con trazabilidad:

\begin{itemize}[leftmargin=*]
    \item ¿Qué productos existen y cuáles se encuentran disponibles?
    \item ¿Cuánto stock queda de cada producto?
    \item ¿Cuánto stock ha ingresado y en qué fecha?
    \item ¿A qué costo se adquirió cada reposición?
    \item ¿Cuánto stock ha salido mediante ventas?
    \item ¿Qué productos están próximos a agotarse?
    \item ¿Cuántas veces se producen quiebres de stock?
\end{itemize}

Adicionalmente, el negocio no dispone de un seguimiento interno detallado que relacione las compras de mercadería, sus costos de adquisición, las ventas realizadas, los precios aplicados y los ingresos obtenidos. Existe información utilizada para obligaciones tributarias, pero no un sistema interno integrado que permita utilizar esos datos para apoyar decisiones operacionales.

\subsubsection{Formulación medible del problema}

El problema central de gestión radica en que el minimarket opera con un 0\% de trazabilidad digital e individual por producto, registrando únicamente el flujo monetario total diario a través del terminal POS sin desglose de artículos. Esta falta de registro unificado obliga a realizar controles de inventario 100\% visuales y manuales, lo que causa quiebres de stock recurrentes en ciclos de reposición rígidos de 15 días, pérdidas por merma no contabilizada y un tiempo perdido estimado de 1 a 2 horas diarias dedicadas exclusivamente a la revisión física de perchas e ingreso manual de datos.

\subsection{Objetivo del SIG propuesto}

El objetivo es desarrollar una aplicación local que permita registrar, controlar y consultar de forma integrada los productos, el inventario, las entradas de stock, los costos de adquisición, las ventas y la información financiera básica del minimarket.

La solución busca reemplazar una administración basada principalmente en observación y experiencia por una gestión sustentada en información registrada. En particular, el sistema deberá apoyar decisiones de reposición, reducir el riesgo de quiebres de stock, conservar información histórica de costos y ventas y entregar indicadores útiles para la operación diaria.

%==============================================
\section{Requerimientos desde cero}
%==============================================

\subsection{Actores y roles}

Los actores identificados hasta el momento son los siguientes:

\begin{longtable}{|p{0.25\textwidth}|p{0.65\textwidth}|}
\hline
\textbf{Actor / sistema} & \textbf{Rol en la solución} \\
\hline
Encargado del minimarket & Usuario principal de la aplicación. Registra productos, entradas de stock y ventas; consulta inventario, costos, historial e indicadores. \\
\hline
Proveedor & Actor externo relacionado con la reposición de mercadería. Entrega productos que posteriormente son registrados como entradas de stock. \\
\hline
Lector de código de barras & Dispositivo utilizado para capturar el identificador de productos durante la carga inicial, búsqueda, reposición y venta. \\
\hline
Google Drive & Servicio externo utilizado únicamente para mantener una copia de respaldo en la nube cuando exista conexión a Internet. \\
\hline
\end{longtable}


\subsection{Alcance}
\begin{itemize} 
\item \textbf{Dentro del Alcance}

El sistema estará orientado principalmente a mejorar la gestión de inventario y ventas del minimarket mediante una aplicación local. La solución permitirá registrar los productos disponibles en el negocio, identificarlos a través de su código de barras y mantener actualizada la cantidad existente de cada uno. Para aquellos productos que no dispongan de código de barras, la aplicación generará automáticamente un identificador propio que permita reconocerlos dentro del sistema.

La aplicación también permitirá registrar las nuevas entradas de mercadería. Cada vez que se incorpore stock de un producto se almacenará la cantidad recibida, el costo asociado a esa adquisición y la fecha del movimiento. De esta forma, además de conocer el inventario actual, se podrá mantener un historial de las distintas reposiciones realizadas.

Las ventas estarán integradas directamente con el inventario. Al registrar una venta, los productos podrán ser identificados mediante el lector de código de barras y, una vez confirmada la operación, la cantidad vendida se descontará automáticamente del stock disponible. Cada venta quedará almacenada junto con los productos involucrados, sus cantidades, el precio utilizado y el monto total de la transacción.

Dentro del sistema también se considerará un apartado de gestión destinado a consultar el estado general del negocio. En este se podrá revisar el inventario disponible, identificar productos con bajo stock, consultar movimientos anteriores y acceder a información básica relacionada con ventas, costos e ingresos.

Cada producto podrá tener definido un nivel mínimo de stock. Cuando la cantidad disponible se acerque o llegue a dicho límite, la aplicación deberá advertir que es necesario realizar una nueva reposición. Este límite buscará considerar el ritmo de venta del producto y el tiempo aproximado que demora la llegada de nueva mercadería.

La aplicación almacenará toda la información localmente en el computador del minimarket y deberá mantener los datos aunque el programa o el computador sean apagados. También se contempla un mecanismo de respaldo de la base de datos, con copias locales y la posibilidad de sincronizar una copia con Google Drive cuando exista conexión a Internet.
\item \textbf{Fuera del Alcance}

En esta primera etapa no se busca desarrollar un sistema contable completo ni reemplazar las herramientas utilizadas para declaraciones tributarias. La información financiera tendrá un carácter principalmente administrativo y permitirá relacionar las ventas realizadas con los costos registrados para obtener información básica sobre el desempeño del negocio.

Tampoco se considera inicialmente el funcionamiento simultáneo desde múltiples computadores o cajas. El sistema estará pensado para funcionar en un único computador Windows y ser utilizado por una persona a la vez.

Algunas funciones adicionales todavía se encuentran pendientes de definición, principalmente el manejo detallado de métodos de pago, cierre de caja, devoluciones, anulaciones de ventas, promociones, productos vendidos por peso y una posible integración con la máquina de pago utilizada por el negocio.

\end{itemize}
\subsection{Requisitos funcionales}

A partir de las necesidades identificadas para el minimarket, se han definido las principales funciones que deberá realizar la aplicación.
\begin{itemize}[leftmargin=*]
\item  \textbf{Registro de productos.}  
El sistema deberá permitir incorporar productos nuevos, registrando como mínimo un identificador, un nombre reconocible y su precio de venta.

\item \textbf{Identificación de productos.}  
Cuando un producto tenga código de barras, este será utilizado como su identificador principal. Si el producto no posee uno, la aplicación deberá generar automáticamente un código propio que no pueda confundirse con un código de barras comercial.

\item \textbf{Ingreso de stock.}  
La aplicación deberá permitir aumentar la cantidad disponible de un producto cuando llegue nueva mercadería al minimarket.

\item \textbf{Registro de costos.}  
Cada nueva entrada de stock deberá almacenar el costo correspondiente a esa adquisición, permitiendo conservar los distintos costos que haya tenido un mismo producto a lo largo del tiempo.

\item \textbf{Historial de movimientos.}  
Las entradas de mercadería deberán quedar registradas junto con su cantidad, costo, fecha y hora, de manera que sea posible revisar posteriormente los movimientos realizados.

\item \textbf{Registro de ventas.}  
El sistema deberá permitir realizar ventas compuestas por uno o más productos, mostrando las cantidades incorporadas, sus precios y el monto total de la operación.

\item \textbf{Uso del lector de código de barras.}  
Durante la creación de productos, ingreso de stock y realización de ventas, la aplicación deberá permitir identificar productos utilizando el lector de código de barras disponible en el local.

\item \textbf{Actualización automática del inventario.}  
Cuando una venta sea confirmada, la cantidad correspondiente deberá descontarse automáticamente del inventario disponible.

\item \textbf{Conservación del historial de ventas.}  
Las ventas deberán permanecer almacenadas con los productos vendidos, cantidades, precios utilizados, fecha, hora y monto total. Si posteriormente cambia el precio actual de un producto, las ventas anteriores deberán conservar el precio que se utilizó originalmente.

\item \textbf{Control de stock mínimo.}  
El sistema deberá permitir establecer un nivel mínimo para cada producto y señalar aquellos que necesiten una próxima reposición.

\item \textbf{Información financiera básica.}  
A partir de los costos registrados y las ventas realizadas, la aplicación deberá permitir consultar información básica sobre ingresos, costos y resultados asociados a la operación del minimarket.

\item \textbf{Indicadores de gestión.}  
La información almacenada deberá utilizarse para calcular indicadores relacionados con inventario y ventas, como productos con bajo stock, ventas diarias, unidades vendidas por producto y precisión del inventario. 


\item \textbf{Respaldo de información.}  
La aplicación deberá permitir generar una copia de seguridad de la base de datos y, cuando exista conexión a Internet, sincronizar una copia con Google Drive.
\end{itemize}

\subsection{Requisitos no funcionales}

Además de las funciones que realizará la aplicación, se establecen algunas condiciones generales que deberá cumplir durante su funcionamiento.
\begin{itemize}[leftmargin=*]
    \item \textbf{Funcionamiento local.}  
    Las operaciones principales deberán ejecutarse directamente en el computador del minimarket mediante \texttt{localhost}, sin necesitar una conexión permanente a Internet.
    
    \item \textbf{Persistencia y respaldo de los datos.}  
    La información deberá mantenerse almacenada aunque se cierre la aplicación, se reinicie Windows o se apague el computador. También deberá existir la posibilidad de generar respaldos para disminuir el riesgo de pérdida de información.
    
    \item \textbf{Compatibilidad con Windows.}  
    El sistema deberá poder instalarse y ejecutarse correctamente en el computador Windows utilizado por el negocio.
    
    \item \textbf{Facilidad de uso.}  
    La interfaz deberá ser clara y sencilla, considerando que estará destinada al uso cotidiano dentro de un minimarket. Las operaciones frecuentes, especialmente la venta y el ingreso de stock, deberán poder realizarse con pocos pasos.
    
    \item \textbf{Compatibilidad con el lector de código de barras.}  
    La aplicación deberá poder recibir e interpretar correctamente los códigos enviados por el lector utilizado en el local.
\end{itemize}

\subsection{Relación de los requerimientos con el problema}

Las funciones definidas buscan responder directamente a las dos principales dificultades observadas en el negocio, levantadas a partir de la entrevista realizada a la encargada del minimarket. Por una parte, el registro de productos, entradas, ventas, stock mínimo e historial permitirá reemplazar el control visual del inventario por información almacenada y actualizada, dado que la encargada confirmó que la detección de bajo stock depende exclusivamente de la observación visual, sin cantidad mínima definida ni registro alguno. Por otra parte, el registro de costos y ventas permitirá contar con información económica básica que actualmente no se encuentra integrada dentro de un mismo sistema, ya que según lo confirmado en la entrevista, no se registra el costo de cada reposición ni la cantidad vendida por producto, conociéndose únicamente el monto total de venta del día.

De esta manera, la aplicación no se limita a registrar información, sino que busca utilizar los datos generados durante la operación diaria para apoyar decisiones de reposición, reducir el riesgo de quedarse sin productos y entregar una visión más clara del movimiento económico del minimarket.

%==============================================
\section{Procesos BPMN}
%==============================================

\subsection{Proceso actual}

Actualmente, la gestión del inventario del minimarket se realiza principalmente de forma manual y visual, no existe un sistema que permita consultar de manera inmediata la cantidad disponible de cada producto ni un registro que relacione directamente las entradas de mercadería con las unidades que posteriormente son vendidas.

De acuerdo con la entrevista realizada a la encargada, la necesidad de reposición se detecta mediante revisión visual diaria de las estanterías, no existe una cantidad mínima definida; la decisión se toma al ojo, apoyándose en la velocidad de venta de cada producto y dejando que la mercadería se agote casi por completo antes de reponer, lo que la propia encargada describe como su práctica habitual.

Detectada la falta de un producto, el pedido se realiza principalmente de forma presencial cuando el proveedor visita el local varios proveedores pasan con alta frecuencia, incluso a diario aunque ocasionalmente se hace por WhatsApp o mediante llamada de un call center. Una vez realizado el pedido, la mercadería llega en un plazo corto. La cantidad mínima a pedir la fija cada proveedor (varía entre 4, 6, 9 o 12 unidades según el producto), y la encargada ajusta la cantidad adicional según la rotación del producto y el espacio disponible en estantería, dado que el local no cuenta con bodega. El pedido queda anotado en un cuaderno personal con el proveedor, el valor a cancelar y la fecha de llegada, pero esta anotación no constituye un registro de inventario.

Al recibir la mercadería, la encargada revisa la factura, las cantidades y las fechas de vencimiento, pero no registra en ningún sistema cuántas unidades ingresaron ni el costo de adquisición, ya que hacerlo durante la atención del local resulta dificil. Los productos se colocan directamente en las estanterías, sin paso por bodega. Las facturas se conservan por una semana y las boletas de venta (efectivo y tarjeta) se revisan semanalmente, pero ninguna de ellas se desagrega por producto.


\begin{figure}[H]
    \centering
    \fbox{
        \parbox[c][6cm][c]{0.88\textwidth}{
            \centering
            \includegraphics[width=1\linewidth]{fotos/Proyecto_sistemas asis.png}
        }
    }
    \caption{Proceso actual de control y reposición de inventario.}
    \label{fig:bpmn-asis}
\end{figure}


\subsection{Proceso propuesto}

Con la incorporación del sistema, la gestión del inventario dejará de depender exclusivamente de la observación del encargado. Cada producto estará registrado dentro de la aplicación y tendrá asociada una cantidad disponible, permitiendo conocer su stock actual y detectar con anticipación cuándo se aproxima a un nivel que requiere reposición.

Cuando llegue nueva mercadería, el producto será identificado mediante su código de barras o mediante el identificador generado por la aplicación. Posteriormente se registrará la cantidad recibida y su costo de adquisición. Una vez confirmada la operación, el sistema aumentará automáticamente el stock disponible y conservará el movimiento dentro del historial.

Durante una venta ocurrirá el proceso inverso. El producto será identificado, se incorporará a la venta y, una vez confirmada la transacción, el sistema almacenará la información correspondiente y descontará automáticamente las unidades vendidas.

La información generada por estos movimientos permitirá que el inventario, las ventas y los costos permanezcan relacionados dentro de una misma aplicación.

\begin{figure}[H]
    \centering
    \fbox{
        \parbox[c][6cm][c]{0.88\textwidth}{
            \centering
            \includegraphics[width=1\linewidth]{fotos/Proyecto_sistemas tobe.png}
        }
    }
    \caption{Proceso propuesto de control y reposición de inventario.}
    \label{fig:bpmn-asis}
\end{figure}


\subsection{Cambios esperados en el proceso}

La principal diferencia entre ambos procesos será la existencia de información actualizada sobre el inventario. Actualmente la necesidad de compra se identifica principalmente observando físicamente los productos. Con el sistema propuesto, esta decisión podrá apoyarse en el stock registrado y en un nivel mínimo definido para cada producto.

Asimismo, las entradas y salidas de mercadería dejarán de ser movimientos aislados. Una entrada de productos aumentará el inventario y conservará su costo de adquisición, mientras que una venta disminuirá automáticamente el stock y generará un registro permanente de la transacción.

Esto permitirá disponer de información histórica que posteriormente podrá utilizarse para analizar ventas, costos, comportamiento del inventario y necesidades de reposición.


\begin{table}[H]
\centering
\small
\begin{tabular}{p{3.2cm}p{5.3cm}p{5.3cm}}
\toprule
\textbf{Aspecto} & \textbf{Proceso actual (As-Is)} & \textbf{Proceso propuesto (To-Be)} \\
\midrule
Detección de la necesidad de compra &
Revisión visual diaria de las estanterías &
Comparación automática entre stock actual y punto de reposición \\
Cantidad a pedir &
Estimación según experiencia, salida habitual y espacio &
Apoyada en historial de ventas y plazo de entrega del proveedor \\
Registro del pedido &
Anotación informal, sin sistema &
Registro estructurado con fecha y proveedor \\
Control de la recepción &
Revisión de factura sin camino ante diferencias &
Verificación explícita de conformidad, con reclamo al proveedor si corresponde \\
Registro de la entrada &
No existe actualización formal del inventario &
Se registra cantidad, costo, fecha y hora de cada entrada \\
Conocimiento del stock disponible &
Se obtiene mirando físicamente la estantería &
Consulta inmediata, actualizada con cada movimiento \\
Registro de ventas &
Solo se conoce el monto total vendido, no el producto ni la cantidad &
Cada venta registra el producto, la cantidad y el precio aplicado \\
Mermas y quiebres &
No se registran &
Se registran como ajustes de inventario \\
\bottomrule
\end{tabular}
\caption{Comparación entre el proceso actual y el proceso propuesto.}
\end{table}



%==============================================
\section{Modelo de datos preliminar}
%==============================================

La aplicación requerirá una base de datos que permita relacionar los productos con los movimientos de inventario y las ventas realizadas. El diseño definitivo dependerá del proceso To-Be, por lo que en esta etapa solamente se identifican las principales estructuras de información necesarias para el funcionamiento del sistema.

El elemento central será el producto. Cada producto deberá contar con un identificador único, un nombre reconocible, un precio de venta actual y la información necesaria para conocer su disponibilidad dentro del inventario. Cuando el producto posea un código de barras, este será utilizado como su identificador principal. En caso contrario, la aplicación generará un identificador propio.

Cada vez que se reciba nueva mercadería se deberá generar un registro independiente de entrada de stock. Este registro permitirá conocer qué producto ingresó, qué cantidad fue recibida, cuál fue su costo de adquisición y cuándo ocurrió el movimiento. Esto es necesario porque un mismo producto puede ser comprado en diferentes momentos y a distintos costos.

Las ventas también deberán mantenerse como registros independientes. Cada venta tendrá una fecha, una hora y un monto total. A su vez, será necesario conservar el detalle de los productos incluidos en ella, indicando qué producto fue vendido, qué cantidad se vendió y cuál fue el precio utilizado en ese momento.

Esta separación permitirá conservar correctamente la información histórica. Por ejemplo, si el precio actual de un producto cambia, las ventas realizadas anteriormente no deberán modificarse.

El modelo deberá permitir que las acciones representadas en el BPMN tengan una correspondencia directa con la información almacenada. Registrar una entrada de mercadería deberá generar un movimiento asociado a un producto; realizar una venta deberá generar una transacción y su correspondiente detalle; y las consultas de inventario deberán obtenerse a partir de esta información.



%==============================================
\section{Arquitectura lógica y stack tentativo}
%==============================================

\subsection{Funcionamiento general}

La aplicación será ejecutada en un único computador ubicado en el minimarket y utilizará una interfaz web accesible desde el navegador mediante \texttt{localhost}. Aunque visualmente funcione como una página web, las operaciones principales se realizarán directamente en el equipo y no dependerán de una conexión permanente a Internet.

El usuario interactuará principalmente mediante el teclado, el mouse y un lector de código de barras. La información ingresada desde la interfaz será procesada por la aplicación local y posteriormente almacenada en una base de datos ubicada en el mismo computador.

De esta manera, el funcionamiento general puede representarse como:

\begin{center}
Usuario y lector de código de barras
$\longrightarrow$
Interfaz web
$\longrightarrow$
Aplicación local
$\longrightarrow$
Base de datos
\end{center}

La conexión a Internet tendrá un propósito secundario. Se utilizará principalmente para mantener una copia de respaldo de la información en la nube, pero una pérdida de conexión no deberá impedir registrar productos, realizar ventas, modificar inventario o consultar información.


\subsection{Organización de la aplicación}

La aplicación se organizará alrededor de tres áreas principales.

La primera estará relacionada con los productos y la adquisición de inventario. Desde este apartado se podrán registrar productos nuevos y posteriormente ingresar nuevas cantidades cuando llegue mercadería. También permitirá conservar los costos asociados a cada reposición.

La segunda estará destinada a las ventas. Su diseño deberá permitir escanear rápidamente los productos, incorporarlos a una transacción y confirmar la venta. Una vez realizada esta acción, la aplicación guardará la operación y actualizará automáticamente el inventario.

La tercera área estará destinada a la gestión del negocio. Desde ella se podrá consultar el inventario actual, identificar productos con bajo stock, revisar ventas anteriores, visualizar costos y acceder a información resumida sobre el funcionamiento del minimarket.

La distribución exacta de estas funciones en pantallas, menús o botones todavía no se encuentra definida, pero esta separación permite establecer la lógica general que tendrá la aplicación.


\subsection{Tecnologías consideradas}

Para la interfaz se considera utilizar HTML, CSS y JavaScript, debido a que la aplicación será utilizada desde un navegador.

Para procesar las operaciones del sistema se considera utilizar Python junto con Flask. Este componente sería responsable de recibir la información de la interfaz, aplicar las reglas de funcionamiento y comunicarse con la base de datos.

Como sistema de almacenamiento se considera utilizar SQLite. Esta alternativa permite mantener la información dentro de un único archivo local y resulta adecuada para una aplicación utilizada en un solo computador y por una persona a la vez.

Por lo tanto, el stack tentativo se compone de:

\begin{center}
HTML, CSS y JavaScript
$\longleftrightarrow$
Python con Flask
$\longleftrightarrow$
SQLite
\end{center}

La combinación de Flask y SQLite asegura una instalación liviana y ejecutable de forma autónoma en el entorno \texttt{localhost} de la máquina del minimarket, eliminando la dependencia de servidores externos y minimizando la complejidad operativa.


\subsection{Persistencia y respaldo}

Toda la información necesaria para el funcionamiento cotidiano deberá almacenarse localmente. Esto permitirá que los productos, el inventario, los costos y las ventas permanezcan disponibles aunque se cierre el navegador, se reinicie la aplicación o se apague el computador.

Además del almacenamiento principal se incorporará un mecanismo de respaldo. La idea es que la aplicación pueda crear una copia de la base de datos y conservarla localmente como medida de seguridad.

Cuando exista una conexión a Internet, esta copia podrá sincronizarse automáticamente con una cuenta de Google Drive utilizada por el negocio. De esta forma existirán dos niveles de protección: el archivo utilizado normalmente en el computador y una copia externa.

El respaldo podrá generarse automáticamente cuando la aplicación se cierre correctamente y también mediante una opción manual que permita al usuario crear una copia cuando lo estime necesario.

Se buscará conservar más de una versión del respaldo para evitar que una eventual copia dañada reemplace inmediatamente a todas las versiones anteriores.

\subsubsection{Política de Respaldo, Conservación y Método de Integración}

Para asegurar la continuidad operativa y la integridad de la información ante fallas físicas del computador, el mecanismo de respaldo se estructurará bajo los siguientes criterios técnicos:

\begin{itemize}
    \item \textbf{Política de Retención y Versión de Respaldo:} El sistema mantendrá una rotación fija de hasta \textbf{5 versiones consecutivas} de la base de datos SQLite (esquema FIFO: \textit{First In, First Out}). Cada archivo de copia incluirá una marca temporal en su nombre (por ejemplo, \texttt{backup\_minimarket\_YYYYMMDD\_HHMM.db}), evitando que un fallo local sobrescriba la totalidad de los datos históricos.
    \item \textbf{Procedimiento de Restauración:} La restauración de la base de datos se realizará mediante una herramienta interna de administración alojada en la interfaz de gestión. Ante un evento de corrupción local, el usuario podrá seleccionar la versión más reciente validada desde un menú desplegable y ejecutar un comando de sobreescritura controlada previo cierre de sesiones activas.
    \item \textbf{Integración con Google Drive:} La sincronización con la nube se implementará utilizando la  Google Drive API v3 en Python mediante autenticación \texttt{OAuth 2.0} (o cuenta de servicio \textit{Service Account}). Los tokens de acceso y credenciales sensibles se resguardarán en variables de entorno locales (archivo \texttt{.env}), garantizando la anonimización de la conexión y el cumplimiento de la norma de protección de datos.
\end{itemize}


%==============================================
\section{Información de gestión e indicadores}
%==============================================

La información almacenada no solo permitirá registrar las operaciones del negocio, sino que también deberá utilizarse para apoyar decisiones cotidianas. Para ello, la aplicación mostrará algunos indicadores construidos a partir del inventario y las ventas.

Uno de los principales será la precisión del inventario. Este indicador permitirá comparar las cantidades registradas por el sistema con las cantidades encontradas físicamente en el minimarket. Su utilización requerirá realizar conteos físicos periódicos que permitan comprobar qué tan bien representa la aplicación la situación real.

También se mostrará qué productos se encuentran bajo su nivel mínimo de stock. La finalidad será advertir que una reposición debe realizarse antes de que el producto llegue completamente a cero.

Relacionado con lo anterior, se podrá observar la cantidad de quiebres de stock, entendidos como situaciones en las que un producto se agota cuando todavía existe demanda.

En la operación diaria será posible consultar el monto total vendido y la cantidad de unidades vendidas de cada producto. Esto permitirá identificar qué productos presentan mayor movimiento y utilizar posteriormente esa información para mejorar las decisiones de compra.

Finalmente, al relacionar los precios de venta con los costos de adquisición, se podrá obtener un margen bruto estimado. Este valor tendrá un propósito exclusivamente administrativo y no representará la utilidad contable definitiva de la empresa.

Como complemento, podría utilizarse una estimación de cobertura de inventario, indicando aproximadamente para cuántos días alcanza el stock disponible según el ritmo reciente de ventas.
Debido a lo observado durante el levantamiento, también será útil controlar las mermas registradas por
período y por producto. Este indicador permitirá identificar pérdidas asociadas a vencimientos, deterioro,
quiebres u otras diferencias que actualmente no pueden medirse.


\subsection{Punto de reposición}

El nivel mínimo de inventario no debería establecerse únicamente de manera arbitraria. La intención es relacionarlo con el consumo habitual del producto y con el tiempo que tarda en llegar una nueva reposición.

Conceptualmente se considera:

\begin{equation}
\text{Punto de reposición}
=
\text{demanda durante el tiempo de reposición}
+
\text{margen de seguridad}
\end{equation}

Por ejemplo, si un determinado producto vende varias unidades diariamente y el proveedor demora algunos días en realizar una nueva entrega, la advertencia debería aparecer antes de que el stock disponible sea insuficiente para cubrir ese período.

\subsubsection{Estimación práctica para el minimarket}

A partir de los datos recolectados en el local, el cálculo automático del punto de reposición se basará en las siguientes reglas operativas:

\begin{itemize}
    \item \textbf{Estimación del consumo diario:} El sistema calculará el consumo promedio diario de cada producto analizando el historial de ventas reales registradas durante los últimos \textbf{14 días} de operación.
    \item \textbf{Tiempo de reposición del proveedor:} Según lo declarado en la entrevista, los distribuidores habituales (como bebidas, lácteos y abarrotes) operan con un tiempo de entrega promedio de \textbf{2 a 3 días} desde que se hace el pedido. Este valor será configurable dentro del sistema según cada proveedor.
    \item \textbf{Margen de seguridad:} Se establecerá un inventario de reserva equivalente al consumo de \textbf{1,5 días} de venta del producto. Esto permite absorber incrementos imprevistos en la demanda o ligeros retrasos en la llegada del camión repartidor sin quedarse sin mercadería.
\end{itemize}

De esta manera, si un producto vende promedio 10 unidades al día y su proveedor tarda 2 días en llegar, el sistema marcará la alerta de reabastecimiento automáticamente cuando el stock caiga a 35 unidades $(10 \times 2 + 15)$.



%==============================================
\section{Gestión financiera básica}
%==============================================

La aplicación utilizará la información de las entradas de stock y de las ventas para entregar una visión financiera básica del funcionamiento del minimarket.

Cada entrada de mercadería conservará su costo de adquisición, mientras que cada venta mantendrá el precio real utilizado en la transacción. Esto permitirá revisar posteriormente cuánto se ha comprado, cuánto se ha vendido y cómo han variado los costos de los productos a lo largo del tiempo.

A partir de estos datos se podrá obtener un margen bruto estimado:

\begin{equation}
\text{Margen bruto estimado}
=
\text{Ingresos por ventas}
-
\text{Costo de la mercadería vendida}
\end{equation}

Este resultado no deberá interpretarse como la utilidad final del negocio, ya que inicialmente no incluirá todos los gastos necesarios para realizar un cálculo contable completo, tales como arriendo, remuneraciones, electricidad, impuestos u otros costos administrativos.

El objetivo de este apartado será entregar información real y ordenada que ayude al encargado a comprender el movimiento económico del negocio y que pueda servir como apoyo para otros procesos administrativos.

\subsubsection{Cálculo del costo cuando los precios de compra cambian}

Como un mismo producto se compra a distintos precios a lo largo del tiempo, el sistema utilizará el método del \textbf{Costo Promedio Ponderado} para saber exactamente cuánto costó la mercadería vendida:

\begin{itemize}
    \item \textbf{Actualización automática:} Cada vez que llegue mercadería con un costo diferente, la aplicación calculará un nuevo costo promedio para ese producto, considerando las unidades que ya quedaban en repisa y las nuevas que están ingresando.
    \item \textbf{Registro al vender:} Al momento de concretar una venta, el sistema tomará ese costo promedio actual y lo guardará congelado junto al detalle de la transacción.
\end{itemize}

De esta forma, aunque el precio de compra a los proveedores suba o baje con los meses, la aplicación sabrá siempre cuánto costó realmente lo que se vendió en cada día y entregará el margen de ganancia real sin enredar a la persona a cargo.

%==============================================
\section{Plan de Entrega 2 y riesgos}
%==============================================

\subsection{Plan de trabajo}

La siguiente etapa comenzará una vez que se encuentre completamente levantado y validado el proceso actual del minimarket. A partir de esta información se podrá cerrar el diseño del proceso propuesto y posteriormente comenzar la implementación de la aplicación.

El trabajo se dividirá en hitos consecutivos, de manera que las funciones principales puedan desarrollarse y probarse antes de incorporar componentes secundarios.

\begin{longtable}{|p{0.02\textwidth}|p{0.3\textwidth}|p{0.28\textwidth}|p{0.12\textwidth}|>{\raggedright\arraybackslash}p{0.12\linewidth}|}
\hline
\textbf{Nº} & \textbf{Hito} & \textbf{Resultado esperado} & \textbf{Responsable}  &Fechas en semanas\\
\hline

1 &
Completar el levantamiento del proceso actual mediante la información entregada por el encargado. &
Proceso actual completamente documentado y validado. &
Charles &21 septiembre\\
\hline

2 &
Finalizar los BPMN As-Is y To-Be. &
Diagramas definitivos del proceso actual y del proceso propuesto. &
Igor &28 de septiembre\\
\hline

3 &
Cerrar el modelo de datos y la estructura de la base de datos. &
Modelo ER definitivo y tablas necesarias para la aplicación. &
Michelle&5 de octubre\\
\hline

4 &
Implementar la estructura inicial de la aplicación local. &
Aplicación ejecutándose mediante localhost en Windows. &
Lorenzo&12 de octubre\\
\hline

5 &
Implementar el registro de productos y la gestión de inventario. &
Productos, entradas de stock, costos e historial funcionando. &
Bruno y Michelle&12-26 de Octubre\\
\hline

6 &
Implementar el sistema de ventas. &
Ventas almacenadas y descuento automático del inventario. &
Charles&26 de Octubre \\
\hline

7 &
Implementar el apartado de gestión e indicadores. &
Consultas de inventario, stock mínimo, ventas, costos y KPI. &
Lorenzo&2 de Noviembre\\
\hline

8 &
Implementar el mecanismo de respaldo. &
Copias locales y sincronización con Google Drive cuando exista conexión. &
Igor &9 de Noviembre\\
\hline

9 &
Realizar pruebas integrales del sistema. &
Validación del funcionamiento conjunto de inventario, ventas, historial y respaldos. &
Igor y Charles&9 de Noviembre\\
\hline

10 &
Preparar documentación y versión final. &
Sistema reproducible y documentación actualizada en el repositorio. &
Todos&16 de Noviembre\\
\hline
\end{longtable}


\subsection{Riesgos del proyecto}

Durante el desarrollo y posterior utilización de la aplicación existen distintos riesgos que deben ser considerados desde la etapa de diseño. Los principales se relacionan con la calidad de los datos iniciales, la consistencia entre ventas e inventario, la disponibilidad de los respaldos y la correcta configuración de los niveles de reposición.

\begin{longtable}{|p{0.27\textwidth}|p{0.27\textwidth}|p{0.36\textwidth}|}
\hline
\textbf{Riesgo} & \textbf{Posible consecuencia} & \textbf{Medida considerada} \\
\hline

Carga inicial incorrecta del inventario. &
El stock mostrado por la aplicación no representa la cantidad existente en el local. &
Realizar un conteo físico durante la carga inicial y verificar las cantidades registradas. \\
\hline

Venta registrada sin actualización correcta del inventario. &
Se generan diferencias entre las ventas realizadas y el stock disponible. &
Registrar la venta y el descuento de stock como una misma operación lógica. \\
\hline

Pérdida o corrupción de la base de datos local. &
Pérdida parcial o total de información histórica. &
Mantener respaldos locales y copias adicionales sincronizadas con Google Drive. \\
\hline

Falla de Internet o de Google Drive. &
No se realiza temporalmente el respaldo en la nube. &
Mantener todas las operaciones principales funcionando localmente y sincronizar posteriormente. \\
\hline

Incompatibilidad del lector de código de barras. &
El proceso de venta o registro de productos se vuelve más lento o no funciona correctamente. &
Realizar pruebas con el lector utilizado realmente por el negocio antes de cerrar la implementación. \\
\hline

Configuración incorrecta del stock mínimo. &
Las alertas pueden aparecer demasiado temprano o demasiado tarde. &
Ajustar los valores utilizando ventas reales, tiempo de reposición y un margen de seguridad. \\
\hline
\end{longtable}

\subsection{Riesgos asociados a datos personales}

La versión definida hasta el momento utilizará principalmente información relacionada con productos, inventario, costos y ventas, por lo que inicialmente no se contempla almacenar información personal de los clientes.

Sin embargo, sí deberán considerarse los datos asociados al acceso a servicios externos, particularmente las credenciales necesarias para realizar los respaldos en Google Drive. También podrían existir datos personales de trabajadores si posteriormente se incorpora un sistema de usuarios.

Por esta razón, el sistema deberá procurar almacenar solamente la información necesaria, limitar el acceso a los respaldos y evitar mantener contraseñas o credenciales directamente en texto plano.

\subsubsection{Análisis de cumplimiento de la Ley 21.719}

Para asegurar que la aplicación cumpla con las exigencias de la nueva Ley de Protección de Datos Personales (Ley 21.719) en Chile, se aplicarán las siguientes medidas concretas:

\begin{itemize}
    \item \textbf{Minimización de datos:} La aplicación no solicitará RUT, nombres, correos ni ningún dato personal de los clientes al momento de realizar las ventas habituales en caja.
    \item \textbf{Resguardo de claves y credenciales:} No se guardarán contraseñas ni claves de acceso escritas directamente dentro del código del programa. Las credenciales de conexión con Google Drive se almacenarán en un archivo local resguardado (usando variables de entorno \texttt{.env}) fuera del repositorio público.
    \item \textbf{Datos de la encargada del local:} Los únicos datos registrados serán el nombre y correo del usuario administrador del minimarket. Esta información se utilizará únicamente para gestionar el acceso local al sistema y no se compartirá con terceros.
    \item \textbf{Uso de datos sintéticos en desarrollo:} Durante las pruebas del sistema en las fases de desarrollo en \texttt{localhost}, se utilizarán únicamente datos de prueba simulados para no exponer información real del negocio ni de sus proveedores.
\end{itemize}


%==============================================
\section{Conclusiones}
%==============================================

La propuesta desarrollada busca solucionar principalmente la falta de información confiable sobre el inventario del minimarket. La aplicación permitirá registrar los productos disponibles, conocer sus existencias, mantener un historial de las nuevas adquisiciones y actualizar automáticamente las cantidades a medida que se realizan ventas.

La integración entre inventario, costos y ventas permitirá que la información generada durante la operación cotidiana pueda utilizarse posteriormente para apoyar decisiones de reposición y para observar de manera básica el comportamiento económico del negocio.

El diseño prioriza el funcionamiento local debido a las características de la PYME y al uso de un único computador. Esto permite que las operaciones principales no dependan de Internet y que la conexión externa sea utilizada solamente como apoyo para mantener respaldos de la información.

Aún se encuentran pendientes algunos elementos relevantes, especialmente el levantamiento completo del proceso actual, los diagramas BPMN definitivos, el modelo entidad--relación y determinadas reglas específicas relacionadas con reposición, costos, métodos de pago y seguridad. Estos aspectos se completarán a medida que se obtenga la información faltante del encargado y se avance hacia la implementación del sistema.


%==============================================
\appendix
\section{Evidencia de la PYME}
%==============================================
En este apartado se muestra una del interior del local, como respaldo de la existencia y el funcionamiento del negocio, además de servir como evidencia de la visita realizada.

\begin{figure}[H]
    \centering
    \includegraphics[width= 0.3\textwidth]{fotos/Evidencia negocio.jpeg}
    \caption{Foto del interior del local. Muestra la distribucion de este.}
    \label{fig:Foto_negocio}
\end{figure}

%==============================================
\section{Entrevista al encargado}
%==============================================
El levantamiento del proceso actual se realizó mediante preguntas dirigidas a la dueña del minimarket. A partir de sus respuestas se confirmó que el inventario se revisa visualmente todos los días, que no existe un registro sistemático de entradas, salidas o mermas y que los pedidos se realizan principalmente de manera presencial cuando pasa cada proveedor.

La dueña indicó que normalmente decide cuánto comprar según la cantidad que queda, la velocidad con que se vende el producto, el espacio disponible y los mínimos exigidos por los proveedores. Los pedidos se realizan aproximadamente cada quince días y la entrega suele ocurrir entre dos y tres días después. Para recordar los compromisos de pago utiliza una libreta donde anota principalmente el proveedor, el monto total y el día esperado de llegada.

En la recepción se revisan factura, cantidades, monto y fechas de vencimiento, y la mercadería se coloca directamente en las estanterías porque el local no cuenta con una bodega. En las ventas no se registra el producto ni la cantidad vendida; se conoce principalmente el total vendido, diferenciando efectivo y tarjeta. También se identificaron mermas por vencimientos, deterioro, huevos quebrados y pequeñas pérdidas de productos de menor tamaño.

Esta información fue utilizada para construir el BPMN As-Is y orientar el diseño del proceso To-Be y los requerimientos funcionales del sistema.
\subsection{Ficha de la entrevista y acreditación}

\begin{itemize}
    \item \textbf{Fecha y entrevistada:} 2 de septiembre de 2026. Franca Valeska Barra Neira (Encargada principal).
    \item \textbf{Empresa y RUT:} Minimarket Alma Neira Muñoz (RUT: 7.327.222-K).
    \item \textbf{Consentimiento:} Acta de Consentimiento Informado firmada (Ley 21.719).
\end{itemize}

\subsection{Puntos clave del levantamiento}

\begin{itemize}
    \item \textbf{Inventario:} Control 100\% visual sin registro digital de stock.
    \item \textbf{Reposición:} Pedidos cada 15 días con entrega en 2 a 3 días.
    \item \textbf{Ventas y costos:} Registro diario en POS sin desglose por producto ni historial de costos.
    \item \textbf{Impacto:} Quiebres de stock y pérdida de 1 a 2 horas diarias en revisiones manuales.
\end{itemize}



%==============================================
\section{Anexos}
%==============================================



\begin{figure}[H]
    \centering
    
    \includegraphics[width=0.3\linewidth]{fotos/Inicio De actividades.pdf}
    \caption{Foto de inicio de actividades}
    \label{fig:placeholder}
\end{figure}


\end{document}